\begin{abstract}
As programming course enrollment continues to increase, instructors are facing growing pressure in assignment grading and personalized feedback delivery. Although traditional manual grading draws on teaching experience, it is often time-consuming, slow in feedback turnaround, and difficult to keep fully consistent across a large class. Existing automatic grading systems can quickly determine whether code passes test cases, but they usually provide only binary results ("correct/incorrect") and limited guidance for improvement. To address these issues, this thesis designs and implements an AI-assisted teaching system for programming courses, with a primary focus on backend services, database modeling, and workflow implementation.

The system establishes a complete "submission-execution-analysis-grading-feedback" loop. Technically, the backend is built with FastAPI, the database uses PostgreSQL, and deployment is containerized for reproducibility. After a student submits code, the system executes it in a controlled environment, records runtime outputs and errors, performs static analysis, and computes final correctness scores based on hidden test cases. To improve grading fairness, the system separates the final score from the static-analysis score: the final score reflects functional correctness, while the static-analysis score evaluates code quality and implementation reasoning.

For feedback generation, the system introduces a Retrieval-Augmented Generation (RAG) pipeline. Course documents, teaching materials, and assignment context are used as retrieval sources to help the language model produce feedback that is better aligned with course content and less likely to include out-of-scope suggestions. On the instructor side, the system supports class management, student management, assignment score review, submission detail review, and report navigation. In addition, the system supports online deployment and public access, allowing teachers and students to log in, submit, and review results off campus.

Integration and testing results show that the system has reached a practical level for course use: the core pipeline runs stably, grading and feedback are functional, and instructor-side features can support small- to medium-scale teaching scenarios. This work demonstrates that, in undergraduate teaching contexts, integrating automatic grading, rule-based scoring, RAG-based feedback, and teaching management into a unified platform is feasible, and it provides an implementation foundation for future work on feedback quality optimization, academic integrity checks, and multi-language extension.


\textbf{Keywords:} Programming education; Automatic grading; RAG; Large language model; Static analysis; Teaching management system

\end{abstract}
